\subsection*{10月23补充练习}
\begin{enumerate}
    \item  (2011$\cdot$湖南卷（文）$\cdot$21)已知平面内一动点 $P$ 到点 ${F}\left( {1,0}\right)$ 的距离与点 $P$ 到 $y$ 轴的距离的差等于1.
    \begin{enumerate}
        \item 求动点 $P$ 的轨迹 $C$ 的方程
        \item 过点 $F$ 作两条斜率存在且互相垂直的直线 ${l}_{1},{l}_{2}$ ,设 ${l}_{1}$ 与轨迹 $C$ 相交于点 $A,B$ , ${l}_{2}$ 与轨迹 $C$ 相交于点 $D,E$ ,求 $\vv{AD} \cdot  \vv{EB}$ 的最小值.
    \end{enumerate}
    \begin{jieda}
        \begin{enumerate}
            \item 设动点 $P$ 的坐标为$(x, y)$,由题意为 $\sqrt{{\left( x - 1\right) }^{2} + {y}^{2}} - \left| x\right|  = 1$. 化简得 ${y}^{2} = {2x} + 2\left| x\right|$ ,
            
            当 $x \geq  0$ 时, ${y}^{2} = {4x}$ ；当 $x < 0$ 时, $y = 0$ .
            
            所以动点 ${P}$ 的轨迹 ${C}$ 的方程为 ${y}^{2} = {4x}\left( {x \geq  0}\right)$ 和 $y = 0\left( {x < 0}\right)$ .
            
            \item 
            设 $A\left( {{x}_{1},{y}_{1}}\right) ,B\left( {{x}_{2},{y}_{2}}\right)$ , $D\left( {{x}_{3},{y}_{3}}\right), E\left( {{x}_{4},{y}_{4}}\right)$

            
            $\vv{AD} \cdot  \vv{EB} = \left( {\vv{AF} + \vv{FD}}\right)  \cdot  \left( {\vv{EF} + \vv{FB}}\right) = \vv{AF} \cdot  \vv{EF} + \vv{AF} \cdot  \vv{FB} + \vv{FD} \cdot  \vv{EF} + \vv{FD} \cdot  \vv{FB} = \left| \vv{AF}\right|  \cdot  \left| \vv{FB}\right|  + \left| \vv{FD}\right|  \cdot  \left| \vv{EF}\right| = \left( {{x}_{1} + 1}\right) \left( {{x}_{2} + 1}\right)  + \left( {{x}_{3} + 1}\right) \left( {{x}_{4} + 1}\right)
            $
            
            由题意知,直线 ${l}_{1}$ 的斜率存在且不为 0,设为 $k$ ,则 ${l}_{1}$ 的方程为 $y = k\left( {x - 1}\right)$ .
            
            由 $\left\{  \begin{array}{l} y = k\left( {x - 1}\right) \\  {y}^{2} = {4x} \end{array}\right.$ ,得 ${k}^{2}{x}^{2} - \left( {2{k}^{2} + 4}\right) x + {k}^{2} = 0$，根据韦达定理有$\left\{\begin{array}{l}
                    {x}_{1} + {x}_{2} = 2 + \dfrac{4}{{k}^{2}}  \\
                    {x}_{1}{x}_{2} = 1
            \end{array}\right.$ 根据两直线垂直可知两直线斜率乘积为$-1$，则可得 $\left\{\begin{array}{l}
                    {x}_{3} + {x}_{4} = 2 + 4k^2  \\
                    {x}_{3}{x}_{4} = 1
            \end{array}\right.$
            
            故$\vv{AD} \cdot  \vv{EB} = 1 + \left( {2 + \dfrac{4}{{k}^{2}}}\right)  + 1 + 1 + \left( {2 + 4{k}^{2}}\right)  + 1 = 8 + 4\left( {{k}^{2} + \dfrac{1}{{k}^{2}}}\right)  \geq  8 + 4 \times  2\sqrt{{k}^{2} \cdot  \dfrac{1}{{k}^{2}}} = {16}
            $
            
            当且仅当 ${k}^{2} = \dfrac{1}{{k}^{2}}$ 即 $k =  \pm  1$ 时, $\vv{AD} \cdot  \vv{EB}$ 取最小值 16.
        \end{enumerate}
    \end{jieda}
    \item (2026$\cdot$普陀区二模$\cdot$20)设 $a > 0,b > 0,m\text{、}t \in  \mathbf{R}$ ,双曲线 $\Gamma  : \dfrac{{x}^{2}}{{a}^{2}} - \dfrac{{y}^{2}}{{b}^{2}} = 1$ 的一条渐近线方程是 $y = 2\sqrt{2}x$ ,点 $P$为 $\Gamma$ 右支上的一点,直线 $l$ 的方程是 $x - {my} - t = 0,O$ 是坐标原点.
    \begin{enumerate}
        \item 若点 $P$ 的坐标为(3,8),求双曲线 $\Gamma$ 的方程
        \item 若直线 $l$ 经过点 $O$，且与 $\Gamma$ 交于 $A$ 、 $B$ 两点,直线 ${PA}$ 、 ${PB}$ 的斜率分别为 ${k}_{1}$ 、 ${k}_{2}$，求 ${k}_{1} \cdot  {k}_{2}$ 的值
        \item 设点 ${F}_{1}$ 是 $\Gamma$ 的左焦点,点 ${A}_{1}$ 、 ${A}_{2}$ 是 $\Gamma$ 的左、右两个顶点,直线 ${A}_{1}P$ 与直线 $x = 1$ 交于点 $M$ ,直线 $l$ 经过点 $P$ 与 $\Gamma$ 的右支交于另外一点 $Q$ ,若 $a = 3$ ,且直线 ${MQ}$ 恒过点 ${A}_{2}$ ,求$\bigtriangleup {F}_{1}{PQ}$ 周长的取值范围.
    \end{enumerate}
    \begin{jieda}
        \begin{enumerate}
            \item $\dfrac{x^2}{a^2} - \dfrac{y^2}{8a^2} = 1$，代入$P(3, 8)$，$\dfrac{9}{a^2} - \dfrac{64}{8a^2} = 1$，解得$a = 1$，故$x^2 - \dfrac{y^2}{8} = 1$.
            \item \ps 背景为第三定义$k_1k_2 = e^2 - 1 = k^2 = 8$

            设$A(x_0, y_0), B(-x_0, -y_0), P(x_p, y_p)$，故$k_1 = \dfrac{y_p - y_0}{x_p - x_0}, k_2 = \dfrac{y_p + y_0}{x_p + x_0}$，$k_1k_2 = \dfrac{y_p^2-y_0^2}{x_p^2-x_0^2}$

            $\dfrac{x^2}{a^2} - \dfrac{y^2}{8a^2} = 1$即$x^2 = a^2 + \dfrac{y^2}{8}$

            $k_1k_2 = \dfrac{y_p^2-y_0^2}{x_p^2-x_0^2} = \dfrac{y_p^2-y_0^2}{\left(a^2+\dfrac{y_p^2}{8}\right) - \left(a^2+\dfrac{y_0^2}{8}\right)} = 8$

            \item \ps 背景为极点极线，$M$在极线（准线$x = 1$）上，则$PQ$过极点($F_2(9, 0)$)，故可以先证明$PQ$过$F_2$，再利用双曲线的第一定义将$\bigtriangleup {F}_{1}{PQ}$ 周长转化到$|PQ|$上.

            根据$a = 3$可知双曲线方程为$\dfrac{x^2}{9} - \dfrac{y^2}{72} = 1$

            设$M(1, n)$，$n \in (-8\sqrt{2}, -4\sqrt{2}) \cup (4\sqrt{2}, 8\sqrt{2})$

            $A_1M: x = \dfrac{4y}{n} - 3$，$A_2M: x = -\dfrac{2y}{n}+3$

            联立$A_1M$和$\dfrac{x^2}{9} - \dfrac{y^2}{72} = 1$可得$\left(\dfrac{4y}{n}-3\right)^2 - \dfrac{y^2}{8} = 9$，一解为$y = 0$，另一解为$P$的纵坐标$y_p = \dfrac{192n}{128-n^2}$，代入直线方程可得$P\left(\dfrac{768}{128-n^2}-3, \dfrac{192n}{128-n^2}\right)$

            联立$A_2M$和$\dfrac{x^2}{9} - \dfrac{y^2}{72} = 1$可得$\left(3-\dfrac{2y}{n}\right)^2 - \dfrac{y^2}{8} = 9$，同理可得$Q\left(\dfrac{-192}{32-n^2}+3, \dfrac{96n}{32-n^2}\right)$

            故$\vv{F_2P} = \left(\dfrac{12n^2-768}{128-n^2}, \dfrac{192n}{128-n^2}\right), \vv{F_2Q} = \left(\dfrac{6n^2-384}{32-n^2}, \dfrac{96n}{32-n^2}\right)$

            故$\vv{F_2P} // \vv{F_2Q}$，故$P,F_2,Q$三点共线

            $C_{\bigtriangleup {F}_{1}{PQ}} = |PQ| + |PQ| + 12$

            设$PQ: x = my + 9$，和$\dfrac{x^2}{9} - \dfrac{y^2}{72} = 1$ 得 $\left( {8{m}^{2} - 1}\right) {y}^{2} + {144my} + 576 = 0$
            
            $\left\{  \begin{array}{l} 8{m}^{2} - 1 \neq  0 \\  \Delta  > 0 \end{array}\right.$，且$P,Q$均在右一支可得$m \in \left(-\dfrac{\sqrt{2}}{4}, \dfrac{\sqrt{2}}{4}\right)$
            
            由韦达定理可得 $\left\{\begin{array}{l}
                 {y}_{1} + {y}_{2} =  - \dfrac{144m}{8{m}^{2} - 1}  \\
                 {y}_{1}{y}_{2} = \dfrac{576}{8{m}^{2} - 1}
            \end{array}\right.$

            $|PQ| = \sqrt{1+m^2} \sqrt{(y_1+y_2)^2 - 4y_1y_2} = \sqrt{1 + {m}^{2}} \cdot  \dfrac{{48}\sqrt{1 + {m}^{2}}}{\left| 8{m}^{2} - 1\right| } = \dfrac{{48}\left( {1 + {m}^{2}}\right) }{1 - 8{m}^{2}} = \dfrac{-6\left( {1 - 8{m}^{2}}\right)  + {54}}{1 - 8{m}^{2}} = \dfrac{54}{1 - 8{m}^{2}} - 6 \in  \lbrack {48}, + \infty )$

            因此$C_{\bigtriangleup {F}_{1}{PQ}} = |PQ| + |PQ| + 12 \in [108, +\infty)$
        \end{enumerate}
    \end{jieda}
\end{enumerate}